\documentclass[11pt,a4paper]{article}
\usepackage{fontspec,amsmath,amssymb,amsthm,microtype}
\setmainfont{lmroman10-regular.otf}[BoldFont=lmroman10-bold.otf,ItalicFont=lmroman10-italic.otf,BoldItalicFont=lmroman10-bolditalic.otf]
\usepackage{xcolor}
\usepackage[most]{tcolorbox}
\usepackage{mathrsfs}
\providecommand{\lt}{<}
\providecommand{\gt}{>}
\setmonofont{DejaVuSansMono.ttf}[BoldFont=DejaVuSansMono-Bold.ttf,ItalicFont=DejaVuSansMono-Oblique.ttf,BoldItalicFont=DejaVuSansMono-BoldOblique.ttf,Scale=0.86]
\newfontfamily\leanfallback{STIXTwoMath-Regular.otf}[Scale=0.86]
\definecolor{leanrule}{gray}{0.62}
\definecolor{leanback}{gray}{0.962}
\definecolor{leantitle}{gray}{0.22}
\definecolor{leanmuted}{gray}{0.38}
\newtcolorbox{leanbox}[1]{enhanced,breakable,colback=leanback,colframe=leanrule,
  boxrule=0.4pt,arc=0.6mm,left=2.4mm,right=2.4mm,top=1.8mm,bottom=1.8mm,
  before skip=0.8em,after skip=1em,fonttitle=\small\bfseries,coltitle=leantitle,
  colbacktitle=leanback,titlerule=0.3pt,title={#1}}
\newlength{\leancharwidth}
\newenvironment{leancode}{\par\smallskip\begingroup\ttfamily\footnotesize\raggedright
  \settowidth{\leancharwidth}{M}\setlength{\parskip}{0pt}\setlength{\parindent}{0pt}
  \frenchspacing}{\par\endgroup\smallskip}
\newcommand{\leanline}[2]{\hangindent=\dimexpr#1\leancharwidth+4\leancharwidth\relax
  \hangafter=1\noindent\hspace*{#1\leancharwidth}#2\par}
\newcommand{\leansource}[1]{\par{\footnotesize\color{leanmuted}#1}\par}
\newcommand{\leanlabel}[1]{\par\smallskip\noindent{\small\bfseries #1}\par}
\usepackage[margin=25mm]{geometry}
\usepackage[unicode,hidelinks]{hyperref}
\hypersetup{pdfauthor={LangevinLean Project}}
\setlength{\emergencystretch}{3em}
\setlength{\parskip}{0.4em}
\setlength{\parindent}{0pt}
\allowdisplaybreaks[2]
\raggedbottom
\title{\texorpdfstring{Adaptation to the genuine Brownian past}{Adaptation to the genuine Brownian past}}
\author{LangevinLean Project}
\date{}
\begin{document}\maketitle
Draft mathematical exposition; these notes may contain mistakes.


The filtration at step \(n\) contains the actual initialization and all additive Brownian input up to \(nh\). Adaptation of the diffusion is obtained from its integral construction; adaptation of the Euler chain follows from its recursion. The actual curvature consequently preserves this same past.



Earlier results used here: \cite[Section~3.54]{LangevinLean2026ULASC}.

\section{\texorpdfstring{Identifying the common initialization and path}{Identifying the common initialization and path}}


\textbf{Claim.} The Euler and global-diffusion initial variables agree, and the actual grid reference equals the endpoint of any canonical solution whose horizon reaches that grid time.



\textbf{Proof.} The diagonal initial coupling has zero squared difference, so its two \(L^2\) variables agree. Uniqueness of the Langevin integral equation identifies the segmented reference with the canonical finite-horizon solution. The additive Brownian input is the actual noise measured from time zero.

\begin{leanbox}{Lean: \texttt{actualBrownianNoise}, \texttt{actualInitial\_\allowbreak{}eq\_\allowbreak{}semigroupInitial}, \texttt{actualReference\_\allowbreak{}eq\_\allowbreak{}solutionOn}}
\leansource{\texttt{actualBrownianNoise} \quad WassersteinMixing/\allowbreak{}ActualAdapted.lean, lines 24–26}
\begin{leancode}
\leanline{0}{def\ actualBrownianNoise\ (t\ :\ ℝ)\ :\ Lp\ (E\ P.dimension)\ 2\ (actualProbability\ P\ V)\ :=}
\leanline{2}{ULA.brownianLangevinNoise\ (actualProbability\ P\ V)\ (CanonicalSemigroup.brownian\ P)}
\leanline{4}{(CanonicalSemigroup.brownian\_preBrownian\ P\ V\ hV)\ 0\ t}
\end{leancode}
\leansource{\texttt{actualInitial\_\allowbreak{}eq\_\allowbreak{}semigroupInitial} \quad WassersteinMixing/\allowbreak{}ActualAdapted.lean, lines 29–40}
\begin{leancode}
\leanline{0}{theorem\ actualInitial\_eq\_semigroupInitial\ :}
\leanline{4}{actualInitial\ P\ V\ hV\ =\ CanonicalSemigroup.initial\ P\ V\ hV}
\end{leancode}
\leansource{\texttt{actualReference\_\allowbreak{}eq\_\allowbreak{}solutionOn} \quad WassersteinMixing/\allowbreak{}ActualAdapted.lean, lines 44–58}
\begin{leancode}
\leanline{0}{theorem\ actualReference\_eq\_solutionOn\ (n\ :\ ℕ)\ (T\ :\ ℝ)\ (hT\ :\ 0\ ≤\ T)}
\leanline{4}{(hn\ :\ (n\ :\ ℝ)\ *\ P.step\ ≤\ T)\ :}
\leanline{4}{actualReference\ P\ V\ hV\ n\ =\ CanonicalSemigroup.solutionOn\ P\ V\ hV\ T\ hT\ ((n\ :\ ℝ)\ *\ P.step)}
\end{leancode}
\leanlabel{How the Lean proof follows the argument}
\texttt{actual\allowbreak{}Initial\_\allowbreak{}eq\_\allowbreak{}semigroup\allowbreak{}Initial} uses the exact diagonal coupling cost. \texttt{actual\allowbreak{}Reference\_\allowbreak{}eq\_\allowbreak{}solution\allowbreak{}On} invokes uniqueness for the actual Brownian integral equation, so no path-identification premise is added.
\end{leanbox}


Earlier results used here: \cite[Section~3.17.2]{LangevinLean2026ULASC}; \cite[Section~3.22.23]{LangevinLean2026ULASC}; \cite[Section~3.25.1]{LangevinLean2026ULASC}; \cite[Section~3.29.4]{LangevinLean2026ULASC}; \cite[Section~3.29.5]{LangevinLean2026ULASC}; \cite[Section~3.44.4]{LangevinLean2026ULASC}; \cite[Section~3.8.17]{LangevinLean2026ULASC}.

\section{\texorpdfstring{Constructing the increasing filtration}{Constructing the increasing filtration}}


\textbf{Construction.} Let \(\mathcal F_n=\sigma(Y_0,B_s:0\le s\le nh)\), interpreted with the actual additive Brownian representatives. Inclusion of time intervals gives \(\mathcal F_n\subseteq\mathcal F_{n+1}\). The constructed reference belongs to \(L^2(\mathcal F_n)\).



\textbf{Proof.} The initial variable and each permitted Brownian evaluation generate the stated sub-sigma-algebra. The Picard solution remains in the closed subspace of random variables measurable with respect to that past; uniqueness identifies it with the reference endpoint.

\begin{leanbox}{Lean: \texttt{actualPast}, \texttt{actualPast\_\allowbreak{}le}, \texttt{actualPast\_\allowbreak{}mono} and 2 more}
\leansource{\texttt{actualPast} \quad WassersteinMixing/\allowbreak{}ActualAdapted.lean, lines 61–63}
\begin{leancode}
\leanline{0}{@[instance\_reducible]\ def\ actualPast\ (n\ :\ ℕ)\ :\ MeasurableSpace\ (ActualSpace\ P)\ :=}
\leanline{2}{ULA.lpLangevinPast\ (actualProbability\ P\ V)\ (CanonicalSemigroup.initial\ P\ V\ hV)}
\leanline{4}{(actualBrownianNoise\ P\ V\ hV)\ ((n\ :\ ℝ)\ *\ P.step)}
\end{leancode}
\leansource{\texttt{actualPast\_\allowbreak{}le} \quad WassersteinMixing/\allowbreak{}ActualAdapted.lean, lines 65–66}
\begin{leancode}
\leanline{0}{theorem\ actualPast\_le\ (n\ :\ ℕ)\ :\ actualPast\ P\ V\ hV\ n\ ≤\ (inferInstance\ :\ MeasurableSpace\ (ActualSpace\ P))}
\end{leancode}
\leansource{\texttt{actualPast\_\allowbreak{}mono} \quad WassersteinMixing/\allowbreak{}ActualAdapted.lean, lines 68–77}
\begin{leancode}
\leanline{0}{theorem\ actualPast\_mono\ :\ Monotone\ (actualPast\ P\ V\ hV)}
\end{leancode}
\leansource{\texttt{actualFiltration} \quad WassersteinMixing/\allowbreak{}ActualAdapted.lean, lines 80–83}
\begin{leancode}
\leanline{0}{def\ actualFiltration\ :\ Filtration\ ℕ\ (inferInstance\ :\ MeasurableSpace\ (ActualSpace\ P))\ where}
\leanline{2}{seq\ :=\ actualPast\ P\ V\ hV}
\leanline{2}{mono'\ :=\ actualPast\_mono\ P\ V\ hV}
\leanline{2}{le'\ :=\ actualPast\_le\ P\ V\ hV}
\end{leancode}
\leansource{\texttt{actualReference\_\allowbreak{}adapted} \quad WassersteinMixing/\allowbreak{}ActualAdapted.lean, lines 86–98}
\begin{leancode}
\leanline{0}{theorem\ actualReference\_adapted\ (n\ :\ ℕ)\ :}
\leanline{4}{AEStronglyMeasurable[actualPast\ P\ V\ hV\ n]\ (actualReference\ P\ V\ hV\ n)\ (actualProbability\ P\ V)}
\end{leancode}
\leanlabel{How the Lean proof follows the argument}
\texttt{actual\allowbreak{}Past} specializes the natural-past construction. \texttt{actual\allowbreak{}Past\_\allowbreak{}mono} proves inclusion of every generating time evaluation and \texttt{actual\allowbreak{}Filtration} packages it. \texttt{actual\allowbreak{}Reference\_\allowbreak{}adapted} uses the proved natural-past theorem for the actual integral solution.
\end{leanbox}


Earlier results used here: \cite[Section~3.25.2]{LangevinLean2026ULASC}; \cite[Section~3.25.3]{LangevinLean2026ULASC}; \cite[Section~3.54.4]{LangevinLean2026ULASC}; \cite[Section~3.54.5]{LangevinLean2026ULASC}; \cite[Section~3.54.8]{LangevinLean2026ULASC}.

\section{\texorpdfstring{Adaptation of innovations, chain and curvature}{Adaptation of innovations, chain and curvature}}


\textbf{Claim.} The normalized innovation is the difference of the actual Brownian inputs at \((n+1)h\) and \(nh\), divided by \(\sqrt{2h}\), hence is measurable by \(\mathcal F_{n+1}\). The Euler state and synchronous error are \(\mathcal F_n\)-measurable, and \(A_n\) maps \(L^2(\mathcal F_n)\) into itself.



\textbf{Proof.} Use the exact increment identity. Induct through the Euler update using continuity of the gradient and the increasing filtration. Subtracting adapted states gives the adapted error. The curvature-preservation theorem then applies because both of its endpoint states are adapted.

\begin{leanbox}{Lean: \texttt{actualNoise\_\allowbreak{}eq\_\allowbreak{}global\_\allowbreak{}increment}, \texttt{actualNoise\_\allowbreak{}adapted}, \texttt{actualChain\_\allowbreak{}adapted} and 2 more}
\leansource{\texttt{actualNoise\_\allowbreak{}eq\_\allowbreak{}global\_\allowbreak{}increment} \quad WassersteinMixing/\allowbreak{}ActualAdapted.lean, lines 102–117}
\begin{leancode}
\leanline{0}{theorem\ actualNoise\_eq\_global\_increment\ (n\ :\ ℕ)\ :}
\leanline{4}{actualNoise\ P\ V\ hV\ n\ =\ (Real.sqrt\ (2\ *\ P.step))⁻¹\ •}
\leanline{6}{(actualBrownianNoise\ P\ V\ hV\ (((n\ +\ 1\ :\ ℕ)\ :\ ℝ)\ *\ P.step)\ -}
\leanline{8}{actualBrownianNoise\ P\ V\ hV\ ((n\ :\ ℝ)\ *\ P.step))}
\end{leancode}
\leansource{\texttt{actualNoise\_\allowbreak{}adapted} \quad WassersteinMixing/\allowbreak{}ActualAdapted.lean, lines 120–140}
\begin{leancode}
\leanline{0}{theorem\ actualNoise\_adapted\ (n\ :\ ℕ)\ :}
\leanline{4}{AEStronglyMeasurable[actualPast\ P\ V\ hV\ (n\ +\ 1)]\ (actualNoise\ P\ V\ hV\ n)\ (actualProbability\ P\ V)}
\end{leancode}
\leansource{\texttt{actualChain\_\allowbreak{}adapted} \quad WassersteinMixing/\allowbreak{}ActualAdapted.lean, lines 144–168}
\begin{leancode}
\leanline{0}{theorem\ actualChain\_adapted\ (n\ :\ ℕ)\ :}
\leanline{4}{AEStronglyMeasurable[actualPast\ P\ V\ hV\ n]\ (actualChain\ P\ V\ hV\ n)\ (actualProbability\ P\ V)}
\end{leancode}
\leansource{\texttt{actualError\_\allowbreak{}adapted} \quad WassersteinMixing/\allowbreak{}ActualAdapted.lean, lines 171–179}
\begin{leancode}
\leanline{0}{theorem\ actualError\_adapted\ (n\ :\ ℕ)\ :}
\leanline{4}{AEStronglyMeasurable[actualPast\ P\ V\ hV\ n]}
\leanline{6}{(actualReference\ P\ V\ hV\ n\ -\ actualChain\ P\ V\ hV\ n)\ (actualProbability\ P\ V)}
\end{leancode}
\leansource{\texttt{actualCurvature\_\allowbreak{}preserves\_\allowbreak{}past} \quad WassersteinMixing/\allowbreak{}ActualAdapted.lean, lines 182–187}
\begin{leancode}
\leanline{0}{theorem\ actualCurvature\_preserves\_past\ (n\ :\ ℕ)}
\leanline{4}{(u\ :\ Lp\ (E\ P.dimension)\ 2\ (actualProbability\ P\ V))}
\leanline{4}{(hu\ :\ AEStronglyMeasurable[actualPast\ P\ V\ hV\ n]\ u\ (actualProbability\ P\ V))\ :}
\leanline{4}{AEStronglyMeasurable[actualPast\ P\ V\ hV\ n]\ (actualCurvature\ P\ V\ hV\ n\ u)\ (actualProbability\ P\ V)}
\end{leancode}
\leanlabel{How the Lean proof follows the argument}
The proof works inside the closed measurable submodule \texttt{lp\allowbreak{}Meas}, whose closure under sums and scalar multiplication handles the recursion. \texttt{actual\allowbreak{}Curvature\_\allowbreak{}preserves\_\allowbreak{}past} supplies the derived adaptation of both actual endpoint states to the actual random-curvature theorem.
\end{leanbox}


Earlier results used here: \cite[Section~3.13.2]{LangevinLean2026ULASC}; \cite[Section~3.25.1]{LangevinLean2026ULASC}; \cite[Section~3.32.13]{LangevinLean2026ULASC}; \cite[Section~3.44.1]{LangevinLean2026ULASC}; \cite[Section~3.54.2]{LangevinLean2026ULASC}; \cite[Section~3.54.6]{LangevinLean2026ULASC}; \cite[Section~3.54.7]{LangevinLean2026ULASC}.

\begin{thebibliography}{ULA-LSI}
\bibitem[ULA-SC]{LangevinLean2026ULASC}
LangevinLean Project. \emph{Unadjusted Langevin under Strong Convexity}. LL-01, 2026. Draft report, content revision d34ead06987c. \url{https://odakyildiz.github.io/LangevinLean/docs/pdf/UnadjustedLangevin/FullProof.pdf}.
\end{thebibliography}

\end{document}
